\documentclass[12pt,a4paper]{article}
\usepackage{geometry}
\geometry{margin=2.5cm}
\usepackage{amsmath,amssymb,amsthm}
\usepackage{graphicx}
\usepackage{hyperref}
\usepackage{booktabs}
\usepackage{natbib}
\usepackage{microtype}
\usepackage{xcolor}

\hypersetup{colorlinks=true,linkcolor=blue,citecolor=blue,urlcolor=blue}

% Theorem environments
\theoremstyle{plain}
\newtheorem{theorem}{Theorem}

% Notation
\newcommand{\Mpl}{M_{\rm Pl}}
\newcommand{\Geff}{G_{\rm eff}}
\newcommand{\GN}{G_N}
\newcommand{\obd}{\omega_{\rm BD}}
\newcommand{\rsc}{\rho_{\rm sc}}
\newcommand{\rhob}{\rho_b}
\newcommand{\Vbary}{V_{\rm bary}}
\newcommand{\Vflat}{V_{\rm flat}}
\newcommand{\Lz}{\Lambda_0}
\newcommand{\CMSTG}{\textsc{cmstg}}

\title{\textbf{Two Structural Obstructions to Galactic Dark Matter\\
               in Non-Minimally Coupled Scalar-Tensor Gravity}}

\author{Christopher R. B. Wilson\\[4pt]
  \small Independent researcher\\
  \small\texttt{C.Isomorph@gmail.com}\\
  \small ORCID: [ORCID-ID]}

\date{2026}

\begin{document}
\maketitle

\begin{abstract}
We prove two structural no-go theorems constraining dark-matter
production in scalar-tensor extensions of general relativity with
non-minimal coupling $F(\Psi) = \tfrac{1}{2} + \Lambda_0\Psi^2$.
\textbf{Theorem~3 ($G_{\rm eff}$ ceiling)}: for any theory with this
coupling form, the Solar-System Cassini PPN bound on the Brans--Dicke
parameter requires $\Lambda_0 < 1.29\times10^{-3}\,\Mpl^{-2}$, giving a
maximum effective gravitational coupling $G_{\rm eff}/G_N \leq 1 +
2\Lambda_0\Psi_0^2 \leq 1.018$.  Flat galactic rotation curves require
$G_{\rm eff}/G_N \approx 3.1$, which demands $\Lambda_0 \geq 0.154\,
\Mpl^{-2}$---incompatible with the Solar-System bound by a factor of
$\sim\!119$.  \textbf{Theorem~4 (shape inversion)}: for any
density-dependent coupling $\beta(\rhob)$ that is monotonically
increasing in baryonic density, the predicted gravitational enhancement
$E_{\rm model}(r) = V_{\rm tot}^2/\Vbary^2$ decreases monotonically
with galactic radius, while flat rotation curves require an enhancement
that increases with radius.  The model and requirement are
anti-correlated by construction; no parameter choice within the
monotone class can produce a positive Pearson coefficient.  Numerical
evaluation on NGC~3198 gives $r_P = -0.81$ and
$\chi^2/\mathrm{dof} = 91.9$.  Together the theorems close the two
principal modification channels in the quadratic-coupling class: $G_{\rm
eff}$ modification and density-dependent fifth-force coupling.  We use
curvature-memory scalar-tensor gravity (\CMSTG) as the worked example
and identify precisely to which theories each theorem applies.
\end{abstract}

\noindent\textbf{Keywords:} scalar-tensor gravity, dark matter, modified gravity,
post-Newtonian constraints, galactic rotation curves

\tableofcontents

%=============================================================================
\section{Introduction}
\label{sec:intro}
%=============================================================================

Scalar-tensor extensions of general relativity (GR) provide one of the
most systematic frameworks for probing modifications to gravity on
astrophysical and cosmological scales~\cite{BransDicke1961,Bergmann1968,
Faraoni2004,Clifton2012}.  The canonical non-minimal coupling
$F(\Psi)\,R$, with $F(\Psi) = \tfrac{1}{2} + \Lambda_0\Psi^2$, is the
leading renormalizable departure from GR: the quadratic term is the
unique dimension-4 operator available for a real scalar, and it subsumes
the Brans--Dicke theory~\cite{BransDicke1961,Damour1992} as a special
case.

Spiral galaxies exhibit flat outer rotation curves that exceed the
Newtonian prediction from observed baryonic matter by factors of 2--5 in
circular velocity, or 4--25 in gravitational potential~\cite{FamaeyMcGaugh2012,
McGaughLelli2016}.  Within scalar-tensor extensions of GR, two broad
mechanism classes address this discrepancy without invoking new matter:
\begin{enumerate}
  \item \emph{Effective coupling modification} ($G_{\rm eff}$ mechanism):
        the background scalar field $\Psi_0$ differs from the field value
        $\Psi_{\rm local}$ inside a galactic halo, raising the effective
        Newton constant $\Geff = G_N\,F(\Psi_0)/F(\Psi_{\rm local})$.
  \item \emph{Density-dependent fifth-force} (chameleon mechanism):
        the scalar couples to baryonic matter with a coupling $\beta(\rhob)$
        that varies with local density, providing an additional acceleration
        in the outer disc where baryonic density is low~\cite{KhouryWeltman2004a,
        KhouryWeltman2004b,HinterbichlerKhoury2010}.
\end{enumerate}

The main results of this paper are:
\setcounter{theorem}{2}
\begin{theorem}[$G_{\rm eff}$ Ceiling]
\label{thm:geff}
For any scalar-tensor theory with coupling $F(\Psi) = \tfrac{1}{2} +
\Lambda_0\Psi^2$, the Solar-System Cassini PPN bound requires $\Lambda_0
< 1.29\times10^{-3}\,\Mpl^{-2}$ and caps $G_{\rm eff}/G_N \leq 1 +
2\Lambda_0\Psi_0^2 \leq 1.018$.  The galactic flat-curve requirement
$G_{\rm eff}/G_N \approx 3.1$ demands $\Lambda_0 \geq 0.154\,\Mpl^{-2}$.
These two constraints have empty intersection; no value of $\Lambda_0$
simultaneously satisfies both.
\end{theorem}

\begin{theorem}[Shape Inversion]
\label{thm:shape}
Let $\beta(\rhob)$ be monotonically increasing in baryonic density
$\rhob$.  Because $\rhob(r)$ decreases monotonically with galactic radius
$r$ outside the disc scale length, the fifth-force enhancement $E_{\rm
model}(r) = 1 + 2\beta(\rhob(r))^2$ is monotonically decreasing with
$r$.  Flat rotation curves require $E_{\rm req}(r) = V_{\rm flat}^2/
\Vbary^2(r)$ to be monotonically increasing with $r$.  The Pearson
correlation between $E_{\rm model}$ and $E_{\rm req}$ is therefore
non-positive for any monotonically increasing $\beta(\rhob)$ and any
galaxy with a falling outer baryonic velocity profile.
\end{theorem}

We prove both theorems, quantify their scope, and illustrate them with the
worked example of curvature-memory scalar-tensor gravity
(\CMSTG)~\cite{WilsonPaperIV}, a covariant scalar-tensor theory with
action
\begin{equation}
  S = \int d^4x\,\sqrt{-g}\Bigl[
    \bigl(\tfrac{1}{2}+\Lambda_0\Psi^2\bigr)R
    - \tfrac{1}{2}(\partial\Psi)^2 - V(\Psi)
  \Bigr] + S_{\rm SM},
  \label{eq:action}
\end{equation}
in units $8\pi G_N = 1$ ($\Mpl = 1$), with canonical locked values
$\Lambda_0 = 0.003\,\Mpl^{-2}$, $\Psi_0 = 2.62\,\Mpl$.

Section~\ref{sec:geff} proves Theorem~\ref{thm:geff} and characterises its
generality. Section~\ref{sec:shape} proves Theorem~\ref{thm:shape} and
identifies what it does and does not rule out.
Section~\ref{sec:conclusions} places the two results in the broader context
of dark-matter model-building.

%=============================================================================
\section{Theorem~\ref{thm:geff}: The \texorpdfstring{$G_{\rm eff}$}{Geff} Ceiling}
\label{sec:geff}
%=============================================================================

\subsection{Effective gravitational coupling and its maximum}
\label{sec:geff_setup}

Consider the action~\eqref{eq:action} with coupling $F(\Psi) = \tfrac{1}{2}
+ \Lambda_0\Psi^2$.  Matter is minimally coupled to the Jordan-frame
metric $g_{\mu\nu}$, so test masses follow geodesics of $g_{\mu\nu}$.
On scales short compared to the inverse scalar mass $m_\Psi^{-1}$, the
effective Newton constant between two test masses depends on the local
scalar field value $\Psi_{\rm local}$:
\begin{equation}
  \frac{\Geff}{\GN} = \frac{F(\Psi_0)}{F(\Psi_{\rm local})}
    = \frac{\tfrac{1}{2}+\Lambda_0\Psi_0^2}{\tfrac{1}{2}+\Lambda_0\Psi_{\rm local}^2},
  \label{eq:geff}
\end{equation}
where $\Psi_0$ is the cosmological background value.  The ratio is
maximised in the limit $\Psi_{\rm local} \to 0$:
\begin{equation}
  \boxed{
  \frac{\Geff}{\GN}\bigg|_{\rm max}
    = \frac{\tfrac{1}{2}+\Lambda_0\Psi_0^2}{\tfrac{1}{2}}
    = 1 + 2\Lambda_0\Psi_0^2.
  }
  \label{eq:ceiling}
\end{equation}
This ceiling depends only on $\Lambda_0$ and $\Psi_0$; it is independent
of the mechanism by which $\Psi_{\rm local}$ is reduced (kernel
accumulation, environmental suppression, etc.).

\subsection{Brans--Dicke correspondence and the Solar-System bound}
\label{sec:geff_pn}

The action~\eqref{eq:action} is equivalent to a Brans--Dicke (BD)
theory~\cite{BransDicke1961} with $\phi_{\rm BD} = 2F(\Psi) = 1 +
2\Lambda_0\Psi^2$.  The BD coupling parameter at the cosmological
background is~\cite{Damour1992,Faraoni2004}
\begin{equation}
  \obd = \frac{F(\Psi_0)}{\bigl(\partial F/\partial\Psi\bigr)^2_{\Psi_0}}
       = \frac{\tfrac{1}{2}+\Lambda_0\Psi_0^2}{4\Lambda_0^2\Psi_0^2}.
  \label{eq:obd}
\end{equation}
For a scalar whose Compton wavelength $m_\Psi^{-1}$ exceeds Solar-System
scales, the post-Newtonian parameter $\gamma$ in BD theory is~\cite{Will2014}
\begin{equation}
  \gamma_{\rm PPN} = \frac{1+\obd}{2+\obd},
  \qquad
  |\gamma_{\rm PPN} - 1| = \frac{1}{2+\obd}.
  \label{eq:gamma}
\end{equation}
The Cassini Shapiro-delay measurement~\cite{Bertotti2003} gives
\begin{equation}
  |\gamma_{\rm PPN} - 1| < 2.3\times10^{-5}
  \qquad \Rightarrow \qquad
  \obd > 43{,}478.
  \label{eq:cassini}
\end{equation}

A key prerequisite for mechanism~(1) (the $\Geff$ modification) to
operate at galactic scales ($r \sim 1$--$50\,{\rm kpc}$) is that the
scalar Compton wavelength must satisfy $m_\Psi^{-1} \gtrsim {\rm few}\,
{\rm kpc}$.  At these Compton wavelengths, the field is also
effectively massless on Solar-System scales ($r_{\rm SS} \sim 1$--$50\,
{\rm AU} \ll {\rm kpc}$), and the PPN formula~\eqref{eq:gamma} with
constraint~\eqref{eq:cassini} applies.

Imposing $\obd > 43{,}478$ in Eq.~\eqref{eq:obd} and solving for
$\Lambda_0$ at $\Psi_0 = 2.62\,\Mpl$ gives
\begin{equation}
  \Lambda_0 < \Lambda_{0,{\rm max}} = 1.29\times10^{-3}\,\Mpl^{-2}.
  \label{eq:lz_max}
\end{equation}
The maximum $\Geff$ ceiling is then
\begin{equation}
  \frac{\Geff}{\GN}\bigg|_{\rm max}^{\rm Cassini}
    = 1 + 2\times(1.29\times10^{-3})\times(2.62)^2
    \approx 1.018.
  \label{eq:ceiling_cassini}
\end{equation}

\subsection{Galactic-scale requirement}
\label{sec:geff_galactic}

For a spiral galaxy with flat outer rotation speed $\Vflat$ and baryonic
circular velocity $\Vbary(r)$, the $\Geff$-modification hypothesis
requires~\cite{FamaeyMcGaugh2012}
\begin{equation}
  \frac{\Geff}{\GN}\bigg|_{r\,\gg\,r_d}
    = \frac{V_{\rm flat}^2}{\Vbary^2(r)}.
  \label{eq:galactic_req}
\end{equation}
For NGC~3198~\cite{Lelli2016} at $r = 30$--$44\,{\rm kpc}$:
$\Vflat \approx 150\,{\rm km\,s^{-1}}$, $\Vbary \approx 85\,
{\rm km\,s^{-1}}$, giving
\begin{equation}
  \frac{\Geff}{\GN}\bigg|_{\rm required}
    = \frac{150^2}{85^2} \approx 3.115.
  \label{eq:geff_req}
\end{equation}
Setting Eq.~\eqref{eq:ceiling} equal to the galactic requirement yields
\begin{equation}
  \Lambda_{0,{\rm req}} = \frac{3.115-1}{2\,\Psi_0^2}
    = \frac{2.115}{2\times(2.62)^2} = 0.154\,\Mpl^{-2}.
  \label{eq:lz_req}
\end{equation}

\subsection{Proof of Theorem~\ref{thm:geff}}
\label{sec:geff_proof}

Equations~\eqref{eq:lz_max} and~\eqref{eq:lz_req} yield:
\begin{align}
  \text{Cassini:}\quad
    & \Lambda_0 < 1.29\times10^{-3}\,\Mpl^{-2},
    \label{eq:cassini_bound}\\[4pt]
  \text{Galactic flat curve:}\quad
    & \Lambda_0 \geq 0.154\,\Mpl^{-2}.
    \label{eq:galactic_bound}
\end{align}
The incompatibility ratio is $0.154/(1.29\times10^{-3}) \approx 119$.
The two allowed regions are disjoint; their intersection is empty.

The severity of the violation at $\Lambda_{0,{\rm req}}$ confirms this is
not a marginal constraint.  At $\Lambda_0 = 0.154\,\Mpl^{-2}$, the BD
parameter (Eq.~\eqref{eq:obd}) drops to $\obd \approx 2.4$, giving
$|\gamma_{\rm PPN} - 1| \approx 0.23$---exceeding the Cassini limit by a
factor of $\sim\!10^4$.  Five additional precision constraints (CMB
temperature spectrum, $f\sigma_8$ at $z < 1$, DESI BAO, UV unitarity, and
recombination-era energy budget) fail simultaneously at
$\Lambda_{0,{\rm req}}$~\cite{WilsonPaperIV}.

Theorem~\ref{thm:geff} is proved. \qquad $\square$

\subsection{Generality and scope}
\label{sec:geff_gen}

\paragraph{Higher-order couplings.}
For $F(\Psi) = \tfrac{1}{2} + \Lambda_n\Psi^n$ the ceiling becomes $1 +
2\Lambda_n\Psi_0^n$ and the BD parameter scales differently with
$\Lambda_n$ and $\Psi_0$, but the same qualitative incompatibility
persists for any renormalizable non-minimal coupling: the Cassini bound
caps $\obd$ from below, which caps $G_{\rm eff}$ from above, while the
galactic requirement sets a minimum.

\paragraph{Brans--Dicke dual statement.}
The incompatibility is equivalent to: the Cassini bound requires
$\obd > 43{,}478$, while NGC~3198 requires $\obd \approx 2.4$.
No value of $\obd$ satisfies both.

\paragraph{Screening and massive scalars.}
A scalar with mass $m_\Psi \gg m_\Psi^{\rm kpc}$ (Compton wavelength
shorter than kpc scales) does not propagate a fifth force at galactic
distances and cannot produce the $\Geff$ modification of
Eq.~\eqref{eq:geff}.  Such theories evade Theorem~\ref{thm:geff}
automatically because they also cannot explain galactic rotation curves
through this mechanism.  The theorem applies precisely in the parameter
regime where the $\Geff$ mechanism is kinematically possible.

\paragraph{What Theorem~\ref{thm:geff} does not rule out.}
The theorem is silent on mechanisms that modify the right-hand side of the
Einstein equations rather than the left-hand side.  Scalar-field dark
matter (ultra-light/fuzzy DM, Q-ball condensates, axion halos) contributes
through $T_{\mu\nu}$ and is not constrained by the PPN bound on $\gamma$.
Theories with non-standard conformal factors or disformal couplings require
a separate PPN analysis.

\paragraph{CMSTG worked example.}
At the locked \CMSTG\ values~\cite{WilsonPaperIV}:
\begin{align}
  \obd &= 2{,}108, & |\gamma_{\rm PPN}-1| &= 4.7\times10^{-4}
  && \text{(Cassini: consistent)},\nonumber\\
  \Geff/\GN|_{\rm max} &= 1.041
  && && \text{($\sim\!3\times$ short of 3.115)}.
\end{align}
A logarithmic sweep of $\Lambda_0$ over eight decades (SIM150)
numerically confirms the ceiling at every value of $\Lambda_0$ consistent
with Solar-System and cosmological constraints.

%=============================================================================
\section{Theorem~\ref{thm:shape}: Shape Inversion for Monotone Couplings}
\label{sec:shape}
%=============================================================================

\subsection{Setup: density-dependent fifth force}
\label{sec:shape_setup}

A fifth-force modification with density-dependent coupling $\beta(\rhob)$
to baryonic matter produces an additional acceleration~\cite{KhouryWeltman2004a,
KhouryWeltman2004b,HinterbichlerKhoury2010}
\begin{equation}
  a_{\rm 5th}(r) = 2\,\beta(\rhob(r))^2\,\GN\,\nabla\Phi_b(r),
  \label{eq:5thforce}
\end{equation}
so that the circular velocity satisfies
\begin{equation}
  V_{\rm tot}^2(r) = \Vbary^2(r)\,\bigl[1 + 2\,\beta(\rhob(r))^2\bigr]
    \equiv \Vbary^2(r)\,E_{\rm model}(r).
  \label{eq:vtot}
\end{equation}
The chameleon and symmetron mechanisms both produce coupling functions
$\beta(\rhob)$ that are monotonically \emph{increasing} with density:
large in the dense inner galaxy and in laboratory environments, small in
the cosmological void where the scalar is screened.

\subsection{Baryon density profile}
\label{sec:shape_density}

Outside the stellar disc scale length $r_d$, the baryonic surface
density falls approximately exponentially, $\Sigma_b(r) \propto e^{-r/r_d}$,
and the mid-plane volume density $\rhob(r)$ tracks this fall.  Gas discs
fall off less steeply, but $\rhob(r)$ is nonetheless monotonically
\emph{decreasing} with $r$ for $r \gtrsim r_d$ in essentially all
rotationally supported late-type galaxies.

\subsection{The shape-inversion argument}
\label{sec:shape_argument}

Let $\beta(\rhob)$ be monotonically increasing.  Since $\rhob(r)$
is monotonically decreasing with $r$, the composition
\begin{equation}
  \beta(r) \equiv \beta(\rhob(r))
\end{equation}
is monotonically \emph{decreasing} with $r$.  Consequently,
\begin{equation}
  E_{\rm model}(r) = 1 + 2\,\beta(r)^2
  \quad \text{is monotonically \emph{decreasing} with } r.
  \label{eq:Emodel_dec}
\end{equation}

For a flat rotation curve with $V_{\rm flat} = {\rm const}$ while
$\Vbary(r) \to 0$ outside the disc, the required enhancement is
\begin{equation}
  E_{\rm req}(r) = \frac{V_{\rm flat}^2}{\Vbary^2(r)},
  \label{eq:Ereq}
\end{equation}
which is monotonically \emph{increasing} with $r$ (since $\Vbary(r)$
decreases).

A monotone-decreasing function and a monotone-increasing function of the
same variable $r$ are anti-correlated.  Their Pearson correlation
coefficient satisfies $r_P \leq 0$, with equality only if one is constant.
No choice of $(\beta_0, \rsc)$ within the monotonically increasing class
can produce $r_P > 0$.

Theorem~\ref{thm:shape} is proved. \qquad $\square$

\subsection{Numerical confirmation on NGC~3198}
\label{sec:shape_num}

Table~\ref{tab:shape} shows $E_{\rm req}$ and $E_{\rm model}$ at four
reference radii for NGC~3198 ($\Vflat = 150\,{\rm km\,s^{-1}}$, $r_d =
2.7\,{\rm kpc}$)~\cite{Lelli2016}, using the tanh profile
$\beta(\rhob) = \beta_0\tanh^2(\rhob/\rsc)$ with self-consistent coupling
$\beta_0 = 1.018$~\cite{WilsonPaperIV}.

\begin{table}[h]
  \centering
  \begin{tabular}{lcccc}
    \toprule
    $r$ [kpc] & 2 & 10 & 20 & 40 \\
    \midrule
    $E_{\rm req}$   & 1.80 & 2.60 & 3.80 & 5.20 \\
    $E_{\rm model}$ & 3.10 & 1.07 & 1.04 & 1.04 \\
    \bottomrule
  \end{tabular}
  \caption{Required and predicted enhancement factors for NGC~3198.
    $E_{\rm model}$ peaks in the inner disc (high density, large $\beta$)
    and saturates near unity in the outer disc; $E_{\rm req}$ rises
    monotonically outward.  The opposite trend in each quantity is the
    structural anti-correlation of Theorem~\ref{thm:shape}.}
  \label{tab:shape}
\end{table}

The Pearson correlation between $E_{\rm req}$ and $E_{\rm model}$ across
all 43 SPARC data points is $r_P = -0.81$, and the best-fit
$\chi^2/{\rm dof} = 91.9$ against an acceptance threshold of $\leq 2.0$.
The optimal transition density $\rsc$ is pinned at its lower boundary,
effectively the uniform limit in which chameleon variation is minimised---the
mechanism is most effective when it varies least, which is a restatement of
the anti-correlation.  These results are reproduced analytically from the
monotonicity argument above.

\subsection{Generality and scope}
\label{sec:shape_gen}

\paragraph{Galaxy and profile independence.}
The proof requires only that $\rhob(r)$ is monotonically decreasing and
$\beta(\rhob)$ is monotonically increasing, both of which hold universally
outside the disc scale length.  The anti-correlation and negative Pearson
coefficient are therefore properties of the entire class of monotonically
increasing $\beta(\rhob)$, not of any specific galaxy or profile parameter.

\paragraph{What Theorem~\ref{thm:shape} does not rule out.}
\begin{enumerate}
  \item \emph{Monotonically decreasing} $\beta(\rhob)$ (``anti-chameleon''):
        a coupling that decreases with density gives $\beta(r)$ increasing
        with $r$, producing $E_{\rm model}(r)$ increasing with $r$.
        Theorem~\ref{thm:shape} is evaded.  However, anti-chameleon profiles
        face an independent obstruction: at cosmological densities $\rhob \sim
        \rho_{\rm crit}$, a decreasing-with-density coupling is large
        (unscreened), generically violating CMB and BAO constraints at order
        unity.  A self-consistent \CMSTG\ analysis constrains the cosmological
        coupling to $\beta_\infty \lesssim 2.4\times10^{-10}$~\cite{WilsonPaperIV},
        which rules out anti-chameleon solutions for that theory.
  \item \emph{Non-monotone} $\beta(\rhob)$: a coupling that peaks at
        intermediate galactic densities and is small at both stellar-core
        and cosmological densities can, in principle, produce $E_{\rm model}$
        peaking in the outer disc.  Theorem~\ref{thm:shape} does not exclude
        this.  The standard thin-shell chameleon always produces monotone
        $\beta(\rhob)$ in its usual operating regime; non-monotone profiles
        require a qualitatively different mechanism architecture.
  \item Mechanisms operating through $T_{\mu\nu}$ sourcing (scalar dark
        matter, condensates, ultra-light fields): these modify the right-hand
        side of the Einstein equations and are not constrained by the
        fifth-force formula~\eqref{eq:5thforce}.
  \item Coupling functions depending on variables other than $\rhob$---for
        example, the tidal field, temperature, or ionization fraction.  For
        these, the argument that $\beta(r)$ is decreasing requires additional
        galaxy-specific input.
\end{enumerate}

\paragraph{CMSTG worked example.}
In \CMSTG, the coupled chameleon coupling is
$\beta(\rhob) = \beta_0\tanh^2(\rhob/\rsc)$ with $\beta_0 = 1.018$ set by
cosmological screening~\cite{WilsonPaperIV}.  The full SPARC rotation-curve
test (SIM152) evaluates this against 43 NGC~3198 data points, confirming
$\chi^2/{\rm dof} = 91.9$ and $r_P = -0.81$.

%=============================================================================
\section{Discussion and Conclusions}
\label{sec:conclusions}
%=============================================================================

We have proved two structural no-go theorems for dark-matter production in
scalar-tensor theories with non-minimal coupling $F(\Psi) = \tfrac{1}{2}
+ \Lambda_0\Psi^2$.

Theorem~\ref{thm:geff} is a \emph{numerical} incompatibility between
two specific observational requirements.  Its conclusion ($\Lambda_0$
incompatible by factor 119) is specific to the combination
$|\gamma_{\rm PPN}-1|_{\rm Cassini}$, $\Vflat$ and $\Vbary$ for NGC~3198,
and $\Psi_0$.  A researcher applying Theorem~\ref{thm:geff} to a different
scalar-tensor theory must re-derive the analog of
Eqs.~\eqref{eq:cassini_bound}--\eqref{eq:galactic_bound} for their own
field value $\Psi_0$ and coupling $\Lambda_0$.  The qualitative
incompatibility is expected to persist for any quadratic non-minimal
coupling.

Theorem~\ref{thm:shape} is a \emph{geometric} (anti-monotone)
incompatibility.  It requires no numerical inputs beyond the fact that
$\rhob(r)$ decreases and $V_{\rm flat}$ is roughly constant in the outer
disc.  The negative Pearson coefficient is guaranteed for any monotonically
increasing $\beta(\rhob)$ and any galaxy with a falling baryonic velocity
profile.  It generalizes to other theoretical frameworks without
re-derivation.

The contrast in character is relevant for future researchers choosing
between the two theorems.  Theorem~\ref{thm:geff} is the sharper tool
when comparing the $G_{\rm eff}$ mechanism against the specific observational
inputs needed for the target theory; Theorem~\ref{thm:shape} is the sharper
tool for ruling out entire classes of monotone couplings without
galaxy-specific computation.

Taken together, the theorems establish that any viable dark-matter
mechanism within the $F(\Psi) = \tfrac{1}{2}+\Lambda_0\Psi^2$ coupling
class must either:
\begin{enumerate}
  \item[(a)] Operate through the energy-momentum $T_{\mu\nu}$ channel
    rather than through a modification of $G_N$ or a density-dependent
    fifth force, or
  \item[(b)] Use a non-monotone coupling $\beta(\rhob)$ with independent
    theoretical motivation beyond the standard chameleon construction.
\end{enumerate}
The theorems do not rule out scalar-tensor theories as hosts for dark
matter; they sharpen which \emph{mechanism} can accomplish this.

For \CMSTG\ specifically, the two theorems (together with the mass-gap
obstruction for the sextic-condensate channel documented
in~\cite{WilsonPaperIV}) exhaust the three available extension channels
for dark-matter production within the locked Phase~1 action.  The
implication is that if \CMSTG\ correctly describes the large-scale
universe, dark matter in the observed universe must come from a sector
structurally distinct from the scalar-tensor framework---standard
candidates (cold dark matter, axions, sterile neutrinos) remain fully
compatible with \CMSTG's cosmological successes.

\section*{Acknowledgements}

The author thanks the open-source scientific Python community (NumPy,
SciPy, Matplotlib) for computational tools used in the numerical
verification.  All simulation code and outputs are publicly archived at
\url{https://github.com/cisomorph/Curvature-Memory_Scalar-Tensor_Gravity}.

\bibliographystyle{unsrtnat}
\bibliography{cmstg_companion_note}

\end{document}
